\section{Билет 47. Теплоемкость. Удельная теплоемкость. Молярная теплоемкость. Теплоемкость идеального газа при изобарическом процессе с выводом. Соотношение Майера.}

\begin{definition}
\textbf{Теплоёмкостью} $C$ тела называется физическая величина, равная количеству теплоты, необходимому для нагревания этого тела на один кельвин:
\begin{equation}
    C = \frac{\delta Q}{dT}. \label{eq:heat_cap_def_47}
\end{equation}
Единица измерения --- Дж/К. Поскольку теплоёмкость пропорциональна массе вещества, вводятся удельные и молярные величины:
\begin{itemize}
    \item \textbf{Удельная теплоёмкость} $c = C/m$ [Дж/(кг·К)] --- теплоёмкость единицы массы;
    \item \textbf{Молярная теплоёмкость} $C_\mu = C/\nu = c M$ [Дж/(моль·К)] --- теплоёмкость одного моля вещества, где $M$ --- молярная масса.
\end{itemize}
\end{definition}

\begin{theorem}[Теплоёмкость идеального газа при постоянном давлении]
Запишем первое начало термодинамики для изобарического процесса ($p = \text{const}$):
\begin{equation}
    \delta Q_p = dU + p \, dV. \label{eq:first_law_isobaric_47}
\end{equation}
Для идеального газа внутренняя энергия зависит только от температуры: $dU = \nu C_V \, dT$. Из уравнения состояния $pV = \nu RT$ при $p = \text{const}$ следует $p \, dV = \nu R \, dT$. Подставляя эти выражения в \eqref{eq:first_law_isobaric_47}, получаем:
\begin{equation}
    \delta Q_p = \nu C_V \, dT + \nu R \, dT = \nu (C_V + R) \, dT.
\end{equation}
По определению молярной теплоёмкости при постоянном давлении $C_p = \frac{1}{\nu} \frac{\delta Q_p}{dT}$. Следовательно:
\begin{equation}
    C_p = C_V + R. \label{eq:Cp_derivation}
\end{equation}
\end{theorem}

\begin{property}[Соотношение Майера]
Уравнение \eqref{eq:Cp_derivation} называется \textbf{соотношением Майера}. Оно показывает, что молярная теплоёмкость идеального газа при постоянном давлении всегда больше теплоёмкости при постоянном объёме на величину универсальной газовой постоянной $R$.
\par\noindent\textbf{Физический смысл:} Разность $C_p - C_V = R$ численно равна работе, совершаемой одним молем идеального газа при его изобарном нагревании на один кельвин. При нагревании при постоянном давлении часть подводимой теплоты расходуется не на увеличение внутренней энергии, а на совершение работы расширения против внешних сил.
\end{property}

\begin{remark}
Используя выражения из классической теории теплоёмкости ($C_V = \frac{i}{2}R$, $C_p = \frac{i+2}{2}R$), соотношение Майера позволяет легко найти \textbf{коэффициент Пуассона} (показатель адиабаты):
\begin{equation}
    \gamma = \frac{C_p}{C_V} = \frac{i+2}{i}.
\end{equation}
Для одноатомных газов ($i=3$): $\gamma \approx 1.67$; для двухатомных ($i=5$): $\gamma = 1.4$; для многоатомных ($i=6$): $\gamma \approx 1.33$. Эти значения хорошо согласуются с экспериментом при обычных температурах.
\end{remark}
